%%
%% The "title" command has an optional parameter,
%% allowing the author to define a "short title" to be used in page headers.
\title{High Throughput Large Integer Multiplier Generation From Optimized FPGA DSP Block Bases}

% \title{ArithV: From Arithmetic to Verilog Circuits on FPGAs}

%%
%% The "author" command and its associated commands are used to define
%% the authors and their affiliations.
%% Of note is the shared affiliation of the first two authors, and the
%% "authornote" and "authornotemark" commands
%% used to denote shared contribution to the research.
% \author{Mahmoud Elshenawy}
% \email{mahmoud.elshenawy@kaust.edu.sa}
% % \orcid{1234-5678-9012}
% \affiliation{%
%   \institution{King Abdullah University of Science and Technology}
%   \city{Thuwal}
%   \country{Saudi Arabia}
% }

% \author{Zhiru Zhang}
% \email{zhiruz@cornell.edu}
% % \orcid{1234-5678-9012}
% \affiliation{%
%   \institution{School of Electrical and\\ Computer Engineering\\}
%   \city{Ithaca, NY}
%   \country{USA}
% }

% \author{Suhaib Fahmy}
% \email{suhaib.fahmy@kaust.edu.sa}
% % \orcid{1234-5678-9012}
% \affiliation{%
%   \institution{King Abdullah University of Science and Technology}
%   \city{Thuwal}
%   \country{Saudi Arabia}
% }



%%
%% By default, the full list of authors will be used in the page
%% headers. Often, this list is too long, and will overlap
%% other information printed in the page headers. This command allows
%% the author to define a more concise list
%% of authors' names for this purpose.
\renewcommand{\shortauthors}{Elshenawy et al.}


% %%
% %% The code below is generated by the tool at http://dl.acm.org/ccs.cfm.
% %% Please copy and paste the code instead of the example below.
% %%
% \begin{CCSXML}
% <ccs2012>
%  <concept>
%   <concept_id>00000000.0000000.0000000</concept_id>
%   <concept_desc>Do Not Use This Code, Generate the Correct Terms for Your Paper</concept_desc>
%   <concept_significance>500</concept_significance>
%  </concept>
%  <concept>
%   <concept_id>00000000.00000000.00000000</concept_id>
%   <concept_desc>Do Not Use This Code, Generate the Correct Terms for Your Paper</concept_desc>
%   <concept_significance>300</concept_significance>
%  </concept>
%  <concept>
%   <concept_id>00000000.00000000.00000000</concept_id>
%   <concept_desc>Do Not Use This Code, Generate the Correct Terms for Your Paper</concept_desc>
%   <concept_significance>100</concept_significance>
%  </concept>
%  <concept>
%   <concept_id>00000000.00000000.00000000</concept_id>
%   <concept_desc>Do Not Use This Code, Generate the Correct Terms for Your Paper</concept_desc>
%   <concept_significance>100</concept_significance>
%  </concept>
% </ccs2012>
% \end{CCSXML}

% \ccsdesc[500]{Do Not Use This Code~Generate the Correct Terms for Your Paper}
% \ccsdesc[300]{Do Not Use This Code~Generate the Correct Terms for Your Paper}
% \ccsdesc{Do Not Use This Code~Generate the Correct Terms for Your Paper}
% \ccsdesc[100]{Do Not Use This Code~Generate the Correct Terms for Your Paper}




%% Keywords. The author(s) should pick words that accurately describe
%% the work being presented. Separate the keywords with commas.
\keywords{
    Arithmetic, 
    Multipliers,
    DSP blocks
}

\maketitle